\begin{theorem}[window-\texorpdfstring{$6$}{6} 退化谱的残数三区间与终端 Zeckendorf 窗]\label{thm:conclusion-window6-residual-threeinterval-terminal-window}
在定理 \ref{thm:conclusion-window6-visible-crt-arithmetic-phase-space} 的规范识别
\[
X_6\cong \ZZ/(21\ZZ)\cong \{0,1,\dots,20\}
\]
下，bin-fold 纤维多重度满足严格分段公式
\[
d_6^{\mathrm{bin}}(w)=
\begin{cases}
4,& 0\le V_6(w)\le 8,\\[1mm]
3,& 9\le V_6(w)\le 12,\\[1mm]
2,& 13\le V_6(w)\le 20.
\end{cases}
\]
并且最小退化壳恰为一个终端 Zeckendorf 长度窗：
\[
S_{6,2}=V_6^{-1}(\{13,\dots,20\})
\cong
13+\{0,1,\dots,7\}
\subset \ZZ/(21\ZZ).
\]
换言之，window-\(6\) 的最小壳不是任意八点子集，而是 \(X_4\) 的规范 Zeckendorf 范围在商半环中的终端平移像。
\end{theorem}
\begin{proof}
由定理 \ref{thm:conclusion-window6-fibonacci-resonance-threshold-lock}，window-\(6\) 的 bin-fold 三相阈值精确对应
\[
V_6\in \{0,\dots,8\},\qquad V_6\in \{9,10,11,12\},\qquad V_6\in \{13,\dots,20\},
\]
其多重度分别为 \(4,3,2\)。这给出分段公式。

另一方面，定理 \ref{thm:conclusion-window6-minimal-shell-rigid-subcover-root-slice} 已证明
\[
S_{6,2}=\iota_4(X_4)=\{x\mathtt{01}:x\in X_4\}.
\]
由 Zeckendorf 范围双射 \(V_4(X_4)=\{0,\dots,7\}\)，再注意后缀 \(\mathtt{01}\) 对 \(V_6\) 的贡献恰为 \(F_7=13\)，遂得
\[
V_6(S_{6,2})=13+\{0,\dots,7\}=\{13,\dots,20\}.
\]
证毕。
\end{proof}

\begin{theorem}[window-\texorpdfstring{$6$}{6} 最小退化壳的 \texorpdfstring{$(5+3)$}{(5+3)} 裂分]\label{thm:conclusion-window6-minimal-shell-five-three-cleavage}
记
\[
X_{6,\mathrm{cyc},2}:=X_6^{\mathrm{cyc}}\cap S_{6,2}.
\]
则存在严格分裂
\[
S_{6,2}=X_{6,\mathrm{cyc},2}\sqcup X_6^{\mathrm{bdry}},
\qquad
|X_{6,\mathrm{cyc},2}|=5,
\qquad
|X_6^{\mathrm{bdry}}|=3.
\]
并且在 \(B_3\) 的 pinned root datum 下，
\[
X_{6,\mathrm{cyc},2}\cong \{-e_1\}\sqcup\{\pm e_2\pm e_3\},
\]
而
\[
X_6^{\mathrm{bdry}}=\{\texttt{100001},\texttt{100101},\texttt{101001}\}
\]
恰对应三重零权。故最小退化壳并非单一根层，而是“五个 cyclic 极小根点 + 三个 boundary 零权点”的刚性拼接。
\end{theorem}
\begin{proof}
由定理 \ref{thm:conclusion-window6-minimal-shell-rigid-subcover-root-slice}，
\[
S_{6,2}\cap X_6^{\mathrm{bdry}}
=
X_6^{\mathrm{bdry}}
=
\{\texttt{100001},\texttt{100101},\texttt{101001}\},
\]
且
\[
S_{6,2}\cap X_6^{\mathrm{cyc}}
=
\{\texttt{000001},\texttt{010001},\texttt{001001},\texttt{000101},\texttt{010101}\}.
\]
因此
\[
S_{6,2}=X_{6,\mathrm{cyc},2}\sqcup X_6^{\mathrm{bdry}},
\qquad
|X_{6,\mathrm{cyc},2}|=5,\quad |X_6^{\mathrm{bdry}}|=3.
\]
同一定理还给出 cyclic 二重点的根字典恰为
\[
\{-e_1\}\sqcup\{\pm e_2\pm e_3\},
\]
而定理 \ref{thm:window6-21-adjoint-weight-multiset} 已把三条 boundary 词识别为三重零权。证毕。
\end{proof}

\begin{theorem}[window-\texorpdfstring{$6$}{6} 有限体积 hidden fiber 的积分同调平凡性]\label{thm:conclusion-window6-hiddenfiber-integral-homology-triviality}
对任意有限指标集 \(\Lambda\) 与任意稳定型配置
\[
w=(w_x)_{x\in \Lambda}\in X_6^\Lambda,
\]
定义 cubical realization
\[
\mathcal F(w):=\prod_{x\in\Lambda} I_{d_6^{\mathrm{bin}}(w_x)}
\subset ([0,1]^2)^\Lambda,
\]
其中 \(I_2,I_3,I_4\subset \{0,1\}^2\) 为 window-\(6\) hidden fiber 的三种仿射 Boolean 方块标准形。则
\[
\widetilde H_k(\mathcal F(w);\ZZ)=0
\qquad (k\ge 0),
\]
并且
\[
\chi(\mathcal F(w))=1.
\]
亦即，任意有限体积下的 hidden fiber 虽可在局部组合上呈现 \(2,3,4\) 三档结构，但其整体积分同调始终退化为可缩型。
\end{theorem}
\begin{proof}
由定理 \ref{thm:xi-time-part9xab-window6-hidden-fiber-affine-boolean-normal-form}，每个单点 hidden fiber 都同构于 Boolean 方块 \(\{0,1\}^2\) 在乘积序下的序理想
\[
I_2,\qquad I_3,\qquad I_4
\]
之一。故 \(\mathcal F(w)\) 恰是这些序理想的有限直积。

另一方面，推论 \ref{cor:xi-time-part9xab-window6-hidden-fiber-global-contractibility} 已证明：任意有限体积下的上述全局 cubical realization 都可缩。于是其全部约化积分同调群在非负次数上皆为零，Euler 特征等于 \(1\)。证毕。
\end{proof}

\begin{theorem}[window-\texorpdfstring{$6$}{6} 仿射纤维的稀疏算术支撑]\label{thm:conclusion-window6-affine-fiber-arithmetic-support}
把
\[
\Omega_6=\FF_2^6
\]
视为仿射空间，并记 \(F_w=(\Fold_6^{\mathrm{bin}})^{-1}(w)\)。则 \(21\) 条纤维中恰有 \(11\) 条为仿射子空间，并且
\[
F_w\ \text{为仿射子空间}
\iff
V_6(w)\in \{0,4,8\}\cup\{13,14,15,16,17,18,19,20\}.
\]
更精确地：
\begin{enumerate}
  \item \(\{13,\dots,20\}\) 上的八条纤维全为仿射 \(1\)-平坦；
  \item \(\{0,4,8\}\) 上的三条四重点纤维全为仿射 \(2\)-平坦；
  \item \(\{9,10,11,12\}\) 上四条纤维全为“仿射平面缺一点”；
  \item 其余 \(\{1,2,3,5,6,7\}\) 上六条四重点纤维全为仿射 \(3\)-平坦的非线性真子集。
\end{enumerate}
故仿射可实现性在 window-\(6\) 中并非按纤维大小单调扩展，而是压缩为
\[
\{0,4,8\}\cup\{13,\dots,20\}
\]
这一稀疏算术支撑。
\end{theorem}
\begin{proof}
定理 \ref{thm:terminal-foldbin6-fiber-affine-geometry} 已证明：window-\(6\) 中恰有 \(8\) 条纤维为 affine\(_1\)-flat、\(3\) 条纤维为 affine\(_2\)-flat，其余 \(10\) 条均非仿射子空间；并且三条 affine\(_2\)-flat 恰对应稳定词
\[
\texttt{000000},\qquad \texttt{101000},\qquad \texttt{000010}.
\]
由 Zeckendorf 值直接计算，
\[
V_6(\texttt{000000})=0,\qquad
V_6(\texttt{101000})=1+3=4,\qquad
V_6(\texttt{000010})=8.
\]

另一方面，由定理 \ref{thm:conclusion-window6-residual-threeinterval-terminal-window}，全部二重点纤维恰落在
\[
V_6\in\{13,\dots,20\},
\]
故这八条纤维正对应全部 affine\(_1\)-flat。

余下四条三重点纤维由定理 \ref{thm:conclusion-window6-residual-threeinterval-terminal-window} 必位于
\[
V_6\in\{9,10,11,12\},
\]
而定理 \ref{thm:terminal-foldbin6-fiber-affine-geometry} 表明它们全部属于 non-affine-plane-minus-one-point 型。最后，除去 \(\{0,4,8\}\) 与 \(\{13,\dots,20\}\) 之后的六个残数
\[
\{1,2,3,5,6,7\}
\]
对应全部剩余四重点纤维，它们由同一定理全部属于 non-affine-subset-of-affine\(_3\)-flat 型。证毕。
\end{proof}
